\section{Discussion}
\label{sec:discussion}

\subsection{Why linearity cannot substitute for accuracy}
\label{sec:why}

The experiment of Section~\ref{sec:trap} is a direct demonstration of an
algebraic fact. Write the reported current as $I_{\mathrm{rep}} = k\,
I_{\mathrm{true}} + c$ for an unknown gain $k$ and offset $c$. Against a
resistive reference $I_{\mathrm{true}} = E/R$, so
$I_{\mathrm{rep}} = (k/R)E + c$. The coefficient of determination of a
least-squares straight-line fit is invariant under affine transformation of the
dependent variable; therefore $R^2$ takes the same value for every $k \neq 0$
and every $c$. A linearity check tests whether the instrument's response is
\emph{proportional} to the stimulus. It says nothing whatever about the
constant of proportionality, which is precisely the calibration.

This is not a subtle point once stated, yet the inference "$R^2 = 0.9999$,
therefore the instrument measures correctly" recurs throughout the low-cost
potentiostat literature and in this instrument's own prior documentation. The
inference is attractive because the premise is genuinely informative---a poor
$R^2$ really would indicate a defective analog chain---and because the failure
mode it misses produces output that looks entirely healthy. In our measurements
the $424\,\%$-in-error condition produced smooth, symmetric, low-noise
voltammograms with a residual structure indistinguishable from the corrected
condition. There is no visual cue.

The practical consequence is that a validation claim built on linearity alone
bounds nothing about absolute current. Any quantity derived from absolute
current---electrode area, rate constants, concentrations, limits of
detection---inherits the unbounded gain error in full. For a comparative or
ratiometric measurement, where a calibration curve is built on the same
instrument and the gain divides out, the error is benign; this distinction is
worth stating explicitly in any report that relies on a linearity check.

\subsection{A defect taxonomy for low-cost instruments}
\label{sec:taxonomy}

The defects in Table~\ref{tab:defects} fall into three groups that call for
different countermeasures.

\emph{Traceability defects} (D1--D3) are mismatches between a software constant
and a hardware fact, or breaks in the path by which a constant reaches the code
that uses it. They are individually trivial and collectively the most dangerous
class, because they alter scale without altering shape. Notably, D1 alone would
have been correctable had D2 not made the calibration resistor unsettable at
run time, and D3 would have masked the correction for the remainder of a
session even so: three independently minor faults compounded into an error that
could not be diagnosed or worked around from the host. The countermeasure is
not more testing but \emph{observability}---an instrument should report the
calibration constants actually in force, which is why check V1 requires
$\hat{R}_{\mathrm{TIA}}$ to be emitted rather than assumed. Adding that one line
of output converted an invisible defect into an obvious one.

\emph{Signal-path and convention defects} (D5--D8) are errors in how the front
end is configured or how its output is interpreted: a programmable-gain enum
read as a gain value, a switch that shunts the measurand, an inverted sign
convention, a power state never established. These are detectable by
differential testing against a known-good path---every one of them was found by
diffing the failing classes against the working sequencer-driven CV engine and
against the vendor reference examples---but not by any end-to-end statistic.

\emph{Robustness defects} (D4) do not corrupt data; they destroy availability.
An unbounded wait on a hardware flag is a single point of failure for the whole
instrument, and it is worth noting that the bounded version we substituted did
not \emph{fix} the impedance path---it converted an indefinite hang into a
reported, diagnosable failure in \SI{50}{\second}. That is the more important
property.

Cutting across all three is an architectural amplifier. The three direct-current
method classes contain verbatim-duplicated configuration and conversion code, so
each of D5--D8 existed in triplicate and had to be corrected three times, and a
fourth, unreachable copy persists in a legacy free-function module. A single
shared conversion and configuration component---the obvious refactoring, and one
this firmware has not yet taken---would have made each defect a single-site fix
and each correction automatically consistent across methods.

\subsection{Per-method validation is not optional}
\label{sec:permethod}

The coverage audit in Section~\ref{sec:coverage} is, we think, the most
transferable result after the linearity argument. On one board, with one cell,
in one session, cyclic voltammetry recovers the reference resistance to better
than half a percent with a replicate spread of order $0.01\,\%$, while
chronoamperometry, square-wave and differential pulse voltammetry on the same
low-power loop and the same electrodes return no cell current at all, and
impedance spectroscopy fails to complete a single frequency point. The methods
differ only in firmware.

It follows that an instrument-level validation claim is not meaningful. A
statement of the form "the instrument was validated against a dummy cell"
should be read as a claim about the one method that was exercised. Where several
methods are implemented---and integrated front ends make implementing many
methods nearly free---each needs its own evidence, and a method that has not
been exercised should be listed as such. We have followed that rule in
Table~\ref{tab:coverage} and in the abstract.

\subsection{Interpretation of the residual error}
\label{sec:residual}

After correction the recovered resistance remains systematically high by
approximately $0.4\,\%$. The component tolerance stack does not account for
this: with $R_{ct}$ specified at $0.1\,\%$ and $R_s$ at a nominal $1\,\%$, the
cell's own uncertainty is about $\pm0.11\,\%$ of $R_{\mathrm{DC}}$. The
remaining $\sim\!0.3\,\%$ is an instrument effect, and the dominant candidate is
the calibration resistor itself: $\hat{R}_{\mathrm{TIA}}$ is proportional to the
\emph{declared} $\Rcal$ by Eq.~(\ref{eq:rtiacal}), and the fitted
\SI{200}{\ohm} part's own tolerance therefore enters the current scale directly
and is not separable from it by any measurement made through that same path.
Characterising the fitted $\Rcal$ against a traceable standard, which we did not
do, would bound this term. Series resistance in the switch matrix and in the
cell wiring adds to the measured value in the same direction and is likewise
not separated here. We therefore report the $0.4\,\%$ as an upper bound on
instrument accuracy under these conditions rather than attributing it.

The stability of the calibration is worth separating from its accuracy. Across
replicates the recovered $\hat{R}_{\mathrm{TIA}}$ varied by only a few parts in
$10^{5}$, so the front end measures its own gain very reproducibly---it was
simply being told to measure it against the wrong number. Reproducibility of
this order is what makes the failure so durable: repeated measurements agree
with each other beautifully, and all of them are wrong together.

\subsection{Limitations}
\label{sec:limitations}

Several limits bound what may be concluded.

\emph{One method is validated, not the instrument.} Only cyclic voltammetry is
supported by quantitative evidence. The remaining five implemented methods are
reported as unvalidated, with three of them positively shown to be
non-functional on this hardware.

\emph{One passive operating point.} The reference is a single passive network of
$\SI{10560}{\ohm}$ at direct current, exercised over $\pm\SI{600}{\milli\volt}$ at
currents up to a few tens of microamperes, with the transimpedance stage at one
nominal setting. Accuracy at other current decades, at other transimpedance
settings, and against faradaic or diffusion-limited responses is not
established. In particular the validated span does not exercise the low-current
range where a biosensing measurement would typically operate.

\emph{The capacitive term was not recovered.} The Randles capacitance predicts a
branch separation of $2\nu C_{\mathrm{eff}}$, which at the highest scan rate
available here remains below the current quantisation step
(Section~\ref{sec:cap}). The cyclic-voltammetry path therefore constrains only
the resistive part of the reference cell, and the protocol's V2 check is a
direct-current check. Recovering $C_{dl}$ would require either a faster ramp
than the firmware's step timing allows or a working impedance path.

\emph{One board, one cell, one operator, one session.} Board-to-board
reproducibility, long-term drift, and environmental sensitivity are not
assessed; the measurement environment was neither controlled nor recorded.
Replicates were acquired consecutively, so temporal drift is not separable from
instrument noise.

\emph{Traceability is to nominal values, not to a standard.} The reference is
the manufacturer's nominal component values and their stated tolerances. No
component was measured against a calibrated instrument, so the entire accuracy
chain rests on those tolerances being honest. This is a deliberate trade---the
protocol is designed to need no reference instrument---but it caps achievable
accuracy at the component tolerance, here of order $0.1\,\%$.
